%! Transformations of Functions — Unit 1, Lesson 1.4 — Lesson Plan
\documentclass[10pt]{article}
\usepackage{precalculus-article}
\usepackage{precalculus-boxes}
\usepackage{graphicx}
\graphicspath{{images/}}
\newcommand{\TallMath}[1]{$\displaystyle #1\rule[-1.4em]{0pt}{3.2em}$}
\newcommand{\CourseName}{Precalculus}
\newcommand{\MeetingLength}{55 minutes}
\newcommand{\UnitNumberName}{Unit 1: Functions and Change \quad}
\newcommand{\LessonNumberName}{Lesson 1.4: Transformations of Functions}

\begin{document}
\begin{center}
    {\Large\bfseries \CourseName \\
    {\small\bfseries \UnitNumberName \LessonNumberName}}
\end{center}

\begin{tcolorbox}[colback=lilac, colframe=plum, boxrule=0.9pt, arc=2mm,
    left=2mm, right=2mm, top=1.5mm, bottom=1.5mm]
    \textbf{Primary Objective:} Students will construct and describe additive
    transformations (vertical and horizontal translations) and multiplicative
    transformations (vertical and horizontal dilations, including reflections)
    of a parent function, using the rule $g(x) = af(x+h)+k$; and will determine
    the domain and range of a transformed function from the parent's domain and
    range. \hfill\textit{\small (CED Topic 1.12 --- Skills 1.C, 3.A)}
\end{tcolorbox}

\renewcommand{\arraystretch}{1.6}
\begin{skillbox}[Priority Ideas \& Skills]{goldbox}
\begin{tabularx}{\linewidth}{@{}X X@{}}
\textbf{AP Skills} &
\textbf{Key Understandings} \\
\midrule
\textbf{1.C} --- Construct new functions, using transformations, that may be
useful in modeling contexts, criteria, or data. &
\begin{itemize}[leftmargin=1.2em,itemsep=2pt,topsep=0pt]
  \item $g(x)=f(x)+k$ is an \textbf{additive} transformation: a vertical
        translation of the graph of $f$ by $k$ units. (1.12.A.1)
  \item $g(x)=f(x+h)$ is an \textbf{additive} transformation: a horizontal
        translation by $-h$ units --- the direction reverses. (1.12.A.2)
\end{itemize} \\
\textbf{3.A} --- Describe the characteristics of a function with varying levels
of precision, depending on the representation available. &
\begin{itemize}[leftmargin=1.2em,itemsep=2pt,topsep=0pt]
  \item $g(x)=af(x)$ is a \textbf{multiplicative} transformation: a vertical
        dilation by a factor of $|a|$; if $a<0$ it also reflects over the
        $x$-axis. (1.12.A.3)
  \item $g(x)=f(bx)$ is a \textbf{multiplicative} transformation: a horizontal
        dilation by a factor of $1/|b|$; if $b<0$ it also reflects over the
        $y$-axis. (1.12.A.4)
\end{itemize} \\
\textbf{Synthesis} --- combine transformations and read off the consequences. &
\begin{itemize}[leftmargin=1.2em,itemsep=2pt,topsep=0pt]
  \item Additive and multiplicative transformations combine:
        $g(x)=af(x+h)+k$. (1.12.A.5)
  \item The domain and range of a transformed function may differ from the
        parent's --- the domain moves with the horizontal changes, the range
        with the vertical ones. (1.12.A.6)
\end{itemize} \\
\end{tabularx}
\end{skillbox}

\begin{skillbox}[{Vocabulary, Concepts, \& Theorems}]{greenbox}
\renewcommand{\arraystretch}{1.4}
\begin{tabularx}{\linewidth}{@{} >{\leavevmode\bfseries\color{plum}}p{0.27\linewidth} X @{}}
Parent function        & The original, untransformed function a whole family of graphs is built from --- today, $f(x)=x^2$. \\
Transformation         & An operation that moves, resizes, or flips the graph of a function without changing its basic shape. \\
Translation (shift)    & A slide: $f(x)+k$ moves the graph up or down; $f(x+h)$ moves it left or right. \\
Dilation (stretch or compress) & A resize: $af(x)$ scales the outputs by $|a|$; $f(bx)$ scales the inputs by $1/|b|$. \\
Reflection             & A flip: $-f(x)$ flips the graph over the $x$-axis; $f(-x)$ flips it over the $y$-axis. \\
Outside vs.\ inside    & A change \textit{outside} $f$ acts on outputs --- vertical, and it does what you expect. A change \textit{inside} $f$ acts on inputs --- horizontal, and it does the opposite. \\
\end{tabularx}
\end{skillbox}

\begin{skillbox}[Activate Prior Knowledge \& Spiral Review]{lilac}
    Please see attached warm-up (spiral review).
    Skills reviewed: evaluating a function at a numerical input with the
    parentheses habit (Lesson 1.1), reading domain and range off a table of
    values (Lesson 1.1), and solving a one-step equation for the input that
    makes an expression zero. Item 3 is the load-bearing one --- ``for what
    $x$ is $x+3 = 0$?'' is exactly the reasoning that explains why $f(x+3)$
    shifts \textit{left}, and it returns verbatim in Part 3.
\end{skillbox}

\begin{skillbox}[Hook]{lilac}
    \textbf{One function, three disguises.} Project the table of $f(x)=x^2$
    beside three unlabeled tables built from it: one with every output 3
    larger, one with the same outputs at inputs 3 smaller, and one with every
    output doubled.
    \begin{enumerate}[itemsep=2pt]
        \item Ask: what did I \textit{do} to $f$ each time? What changed ---
              the inputs or the outputs?
        \item Push on the middle table: its output list is \textit{identical}
              to $f$'s. Nothing happened to the outputs at all --- so what
              moved?
    \end{enumerate}
    \textit{Discussion prompt:} every graph today is that same parabola
    wearing a disguise. Today's job is to read the rule and predict the
    disguise --- and then to reverse it.
\end{skillbox}

\begin{skillbox}[Lesson]{lilac}
\begin{multicols}{2}

\textbf{Part 1: Parent Functions and the Outside/Inside Rule} (8 min)

\begin{itemize}
  \item Name the \textbf{parent function} $f(x)=x^2$ and its five reference
        points $(-2,4)$, $(-1,1)$, $(0,0)$, $(1,1)$, $(2,4)$. Everything
        today is that parabola, relocated or resized.
  \item Give the organizing idea once, in these words, and repeat it all
        lesson: \textbf{outside the function $\Rightarrow$ vertical, and it
        does what you expect; inside the function $\Rightarrow$ horizontal,
        and it does the opposite.}
  \item Before computing anything with a new rule, sort it aloud: is the
        change outside $f$ or inside $f$?
\end{itemize}

\textbf{Part 2: Vertical Translation, $g(x)=f(x)+k$} (8 min)

\begin{itemize}
  \item Fill the notes table for $g(x)=f(x)+3$: every output rises by 3,
        every input untouched. Read the pre-drawn pair of parabolas.
  \item $k>0$ up, $k<0$ down. Domain unchanged; the range slides by $k$.
\end{itemize}

\textbf{Part 3: Horizontal Translation, $g(x)=f(x+h)$} (12 min)

\begin{itemize}
  \item The hard one --- budget the time. Fill the table for $g(x)=f(x+3)$
        and let students discover that the output \textit{list} is unchanged;
        only the addresses moved.
  \item Explain the reversal with the warm-up's question, not with a slogan:
        the parent's vertex sits where its input is $0$, so $g$'s vertex sits
        where $x+3=0$, that is, at $x=-3$ --- \textbf{left} 3.
\end{itemize}

\columnbreak

\textbf{Part 4: Dilations and Reflections} (14 min)

\begin{itemize}
  \item $g(x)=2f(x)$: outputs doubled --- a vertical stretch by 2. Then
        $h(x)=-f(x)$: every output changes sign --- a flip over the
        $x$-axis. All three curves are on the pre-drawn figure.
  \item $|a|>1$ stretches, $0<|a|<1$ compresses. Say ``stretch away from the
        $x$-axis,'' not ``makes it bigger.''
  \item Horizontal dilation at exposure depth: $g(x)=f(2x)$ reaches the same
        heights at \textit{half} the $x$ --- inside, so horizontal, and the
        opposite of what the 2 suggests: a compression by a factor of
        $1/2$.
\end{itemize}

\textbf{Part 5: Combining Them --- and Domain \& Range} (8 min)

\begin{itemize}
  \item Decode $g(x)=-2f(x+1)-3$ one parameter at a time: left 1, stretch by
        2, flip over the $x$-axis, down 3.
  \item Domain and range: if $f$ has domain $[-2,4]$ and range $[0,9]$, the
        domain shifts left 1 to $[-3,3]$; the range multiplies by $-2$ to
        $[-18,0]$ (the endpoints swap order!) and then drops 3 to
        $[-21,-3]$.
  \item State the rule of thumb explicitly: \textbf{horizontal changes move
        the domain, vertical changes move the range.}
\end{itemize}
\end{multicols}
\end{skillbox}

\begin{skillbox}[Active Monitoring]{redbox}
While students work on guided notes and practice, circulate and check:
\begin{itemize}
  \item \textbf{Common error:} shifting $f(x+3)$ to the \textit{right}.
        Cold-call: ``At what $x$ does $g$ do what $f$ did at $0$? Solve
        $x+3=0$.'' Do not accept ``because it's backwards.''
  \item \textbf{Common error:} in the table for $g(x)=2f(x)$, doubling the
        $x$ column. Ask: ``Is the 2 inside or outside $f$? So which column
        does it touch?''
  \item \textbf{Common error:} a range under a negative factor written in the
        old order, e.g.\ $[0,-18]$. ``Smaller number first --- which of those
        two \textit{is} smaller?''
  \item \textbf{Common error:} changing the domain when the transformation
        was vertical. Return to the rule of thumb: horizontal moves the
        domain, vertical moves the range.
\end{itemize}
\end{skillbox}

\begin{skillbox}[Group Work \& Differentiation]{redbox}
Students work on the \textbf{Group Activity: The Skate Park Ramp} in leveled
tiers. Every tier uses the same table of $f$ values and the same pre-drawn
parabola --- no tier draws anything.

\begin{itemize}
  \item \textbf{Tier R (Remediate):} Fill the output columns for
        $g(x)=f(x)+3$ (the ramp raised onto a platform) and $h(x)=f(x)-1$
        (the ramp sunk into a pit), name each direction, then answer the
        inside/outside sorting question.
  \item \textbf{Tier A (Apply):} Build the table for $g(x)=2f(x-1)$ in two
        recorded steps, describe both moves in words, locate the new lowest
        point $(1,0)$, and state the domain $[-1,3]$ and range $[0,8]$.
  \item \textbf{Tier E (Extend):} Reverse-engineer two rules from tables
        ($k(x)=-2f(x)$ and $m(x)=f(x-2)+1$), then compare $f(2x)$ against
        $2f(x)$ at supplied inputs to show the two are genuinely different
        transformations.
\end{itemize}
\end{skillbox}

\begin{skillbox}[Individual Work \& Assessment]{redbox}
\textbf{Exit Ticket} (last 5--7 minutes, individual, collected):
\begin{enumerate}
  \item Decompose $g(x)=f(x+3)-2$ into its horizontal move (with direction)
        and its vertical move (with direction).
  \item Complete a table for $g(x)=-f(x)$ from given values of $f$.
  \item Given that the domain of $f$ is $[-2,6]$, state the domain of
        $g(x)=f(x-4)$ (worked space provided).
\end{enumerate}
\textbf{Assessment note:} Item 1 is the diagnostic --- a ``right 3'' there is
the lesson's signature error and means Part 3 needs a reteach in tomorrow's
warm-up window. Item 3 misses split two ways: $[-6,2]$ is the same reversal
as item 1, while a student who changed the \textit{range} instead confused
which quantity a horizontal move touches.
\end{skillbox}

\begin{skillbox}[Reinforcement \& Extension]{goldbox}
\textbf{Homework:} 5 problems covering description of four single
transformations, completing transformation tables for $f(x)+4$ and $-f(x)$,
domain and range under a combined transformation, reverse-engineering a rule
from two tables, and one AP-style multiple-choice range item.

\textbf{Extension (optional):} Using supplied values of $f(x)=x^2$, compare
$2f(x)$ with $f(2x)$ at the same inputs. The two agree nowhere except the
origin --- outside and inside are genuinely different operations, not two
spellings of one.

\textbf{Preview:} Next lesson (1.5, Composition of Functions) keeps feeding
one function's output into another expression. $f(x+3)$ was already a
composition in disguise; 1.5 gives the idea a name and a notation.
\end{skillbox}

\begin{teachernote}[Warm-Up]
Five minutes, silent start, no calculators. Item 3 is the one that pays off:
``for what $x$ is $x+3=0$?'' is a throwaway algebra question here and the
entire explanation of the horizontal reversal in Part 3 --- so when you get
there, make a point of saying they already answered it. Item 1 keeps the
parentheses habit alive for the tables all lesson; item 2's table read is the
domain/range language Part 5 needs.
\end{teachernote}

\begin{teachernote}[Guided Notes]
Say the outside/inside rule in Part 1 and then repeat it verbatim at every new
rule --- it is the one sentence a struggling student should leave with.
Part 3 is where the time goes; resist explaining the reversal with ``it's
backwards'' and use the $x+3=0$ argument every time. Part 4's horizontal
dilation is exposure only (CED 1.12.A.4): the table showing $f(2x)$ hitting
the same heights at half the $x$ is enough, and no student needs to produce a
horizontal dilation from scratch today. If time runs short, compress that
table into a single read --- do not cut Part 5, which the exit ticket's item
3 depends on.
\end{teachernote}

\begin{teachernote}[Group Activity]
All three tiers read the same ramp table and the same parabola, so groups can
be regrouped mid-activity without new setup. Tier R's last question --- did
the rule change the input column or the output column? --- is that tier's
whole point; make them answer it out loud before moving on. Tier A's two-step
table (shift first, then stretch) is worth insisting on even for students who
can do it in one jump: the recorded intermediate column is what makes their
error visible when it comes. Tier E's $f(2x)$ versus $2f(x)$ comparison is
the best whole-class share of the day if you have three minutes at the end.
\end{teachernote}

\begin{teachernote}[Exit Ticket]
Collected, completion credit only. Sort by item 1: ``left 3, down 2'' is the
target, and the ``right 3'' pile is your reteach group. Accept any wording
that names both a direction and a number of units; do not accept ``moves 3''
with no direction. On item 3, a student who wrote $[2,10]$ with nothing in
the work space still gets credit --- the inequality chain is a support, not a
requirement.
\end{teachernote}

\begin{teachernote}[Homework]
Problem 4 (reverse-engineering $p$ from a table) is the one that predicts
success in 1.7's modeling work; scan it first. On problem 5, distractor A is
the range computed with $a=+2$ instead of $-2$, C drops the $+1$, and D
forgets to reorder the endpoints after multiplying by a negative --- each is
worth thirty seconds of autopsy tomorrow. The extension is the only place
horizontal dilation is practiced outside Tier E; if the class never reached
that tier, assign it explicitly.
\end{teachernote}
\end{document}
